\documentclass[11pt]{amsart}
\usepackage{amsmath,amssymb,amsthm,mathtools}
\usepackage[margin=1in]{geometry}
\usepackage{enumerate}
\usepackage{hyperref}

\newtheorem{theorem}{Theorem}[section]
\newtheorem{proposition}[theorem]{Proposition}
\newtheorem{lemma}[theorem]{Lemma}
\newtheorem{corollary}[theorem]{Corollary}
\theoremstyle{definition}
\newtheorem{definition}[theorem]{Definition}
\theoremstyle{remark}
\newtheorem{remark}[theorem]{Remark}

\newcommand{\RH}{\mathrm{RH}}
\renewcommand{\Re}{\operatorname{Re}}
\renewcommand{\Im}{\operatorname{Im}}
\newcommand{\half}{\tfrac12}
\newcommand{\dnu}{d\nu}
\newcommand{\dmfull}{dm_{\mathrm{full}}}
\newcommand{\dmon}{dm_{\mathrm{full,on}}}
\newcommand{\zon}{\zeta_{\mathrm{on}}}
\newcommand{\Tlam}{T_\lambda}
\newcommand{\Wlam}{W_\lambda}

\title[A Trace--Measure Criterion]{A trace and measure criterion for the Riemann Hypothesis\\
from the zeta spectral triple}
\author{David Alejandro Trejo Pizzo}
\thanks{Independent Researcher. \texttt{dtrejopizzo@gmail.com}}


\begin{document}
\nocite{bgConnes99,bgCC21}
\begin{abstract}
The zeta spectral triple attaches to $\zeta$ two spectral measures: the full measure
$\dmfull$, built from $|\zeta(\half+is)|^2$, and the on-critical measure $\dmon$, built from
the Hadamard product with every zero forced onto the critical line. We study their difference,
the signed measure $\dnu=\dmfull-\dmon$, and prove that it is the exact obstruction to the
Riemann Hypothesis. Our central unconditional result is a \emph{dominance theorem}: the
density $R(s)=|\zeta(\half+is)/\zon(\half+is)|^2$ satisfies $R\ge1$ for all real $s$, with
equality if and only if every nontrivial zero lies on the line; hence $\dnu\ge0$
unconditionally, and $\dnu=0\iff\RH$. From this single non-negative measure we derive a
synthesis: eight a priori different conditions on $\zeta$ --- a zero measure, the vanishing of
a Nevanlinna functional, a Jensen excess, a relative entropy, a conditional-positive trace,
and the others --- are pairwise equivalent, all reducing to $\dnu=0$. We give the
quantitative lower bound that an off-critical zero forces on the trace $\Tlam=\int\Wlam\,\dnu$
(second order in the off-line displacement), and we are careful about the one direction that is
not unconditional: recovering $\dnu$ from the trace family $\{\Tlam\}$ requires the totality of
the Abel kernels in $L^2(\dnu)$, which we state as an explicit completeness hypothesis and
connect to the moment-rigidity criterion for $\RH$ proved elsewhere; we do not use, and we
explicitly avoid, any pointwise positivity of the Abel kernel as a quadratic form. We frame the
obstruction through the function-field analogue, where the same measure vanishes by Weil's
three structural pillars (a polarized cohomology, a Frobenius, a Lefschetz pairing), and
isolate the missing pillar for $\zeta$ as the Hadamard barrier: the Hadamard product delivers
$\dnu\ge0$ but not the completeness or the positivity that would force $\dnu=0$. Nothing here
assumes or proves the Riemann Hypothesis; the contribution is a clean non-negative invariant
whose vanishing is $\RH$, together with a precise account of what its vanishing would require.
\end{abstract}
\maketitle

\tableofcontents
\newpage

\section{Introduction}

The zeta spectral triple represents $\zeta$ through a self-adjoint Jacobi operator whose
spectral data encode the nontrivial zeros. Two measures arise naturally. The \emph{full}
spectral measure $\dmfull$ carries the actual modulus $|\zeta(\half+is)|^2$; the
\emph{on-critical} measure $\dmon$ carries the modulus of the function $\zon$ obtained from the
Hadamard product of $\zeta$ by relocating every zero onto the critical line. If the Riemann
Hypothesis ($\RH$) holds, $\zeta=\zon$ and the two coincide; in general they differ, and the
difference
\[
\dnu=\dmfull-\dmon
\]
is a signed measure on $\mathbb R$ that vanishes exactly when $\RH$ holds. This paper is a study
of $\dnu$ as an invariant in its own right.

The main unconditional fact is that $\dnu$ has a sign. Writing $R(s)=|\zeta(\half+is)/\zon
(\half+is)|^2$ for the density of $\dmfull$ relative to $\dmon$, we prove $R(s)\ge1$ for all
real $s$ (Theorem~\ref{thm:dominance}), with equality iff there are no off-critical zeros; thus
$\dnu\ge0$ unconditionally, and $\dnu=0\iff\RH$ (Corollary~\ref{cor:dnu}). The mechanism is
elementary and exact: each off-critical zero, together with the three partners forced on it by
the functional equation and conjugation, contributes to $R$ a factor $\ge1$ that degenerates to
$1$ precisely when the off-line displacement vanishes.

This non-negative measure organizes a number of previously separate criteria. We prove a
synthesis theorem (Theorem~\ref{thm:eight}): eight conditions --- among them the vanishing of a
Nevanlinna transform of $\dnu$, the vanishing of a Jensen excess on a disk model of $\zeta$, the
vanishing of a relative entropy, and the vanishing of the conditional-positive trace family
$\{\Tlam\}$ --- are all equivalent to $\dnu=0$, hence to $\RH$. We then turn to the trace
$\Tlam=\int\Wlam\,\dnu$ with $\Wlam$ the positive Abel kernel. Because $\dnu\ge0$ and $\Wlam>0$,
each $\Tlam\ge0$ \emph{without cancellation}, and a single off-critical zero forces a strictly
positive lower bound that is second order in its displacement (Proposition~\ref{prop:lower}).
The one delicate point --- which we treat with care --- is the reverse implication: that
$\Tlam=0$ for all $\lambda$ forces $\dnu=0$. This requires the totality of $\{\Wlam\}$ in
$L^2(\dnu)$; we isolate it as an explicit completeness hypothesis and relate it to the
moment-rigidity criterion for $\RH$ established through orthogonal polynomials on the line
\cite{MomentRigidity}. We do \emph{not} appeal to any pointwise sign of $\Wlam$ as a quadratic
form; the only positivity used is that of $\dnu$ itself and of the scalar test functions.

Finally we locate the obstruction structurally. In the function-field analogue the same measure
vanishes \emph{for a proven reason}: Weil's three pillars --- a polarized cohomology with a
positive intersection form, a Frobenius whose eigenvalues are the zeros, and a Lefschetz/Hodge
pairing forcing them onto the critical circle. For $\zeta$ the first two have candidate
analogues but the polarization is absent; what remains is the \emph{Hadamard barrier}: the
Hadamard product supplies $\dnu\ge0$ but supplies neither the completeness nor the positive
intersection form that would force $\dnu=0$. We make this comparison precise
(\S\ref{sec:weil}--\S\ref{sec:barrier}). Nothing in the paper assumes or proves $\RH$.

\section{The two measures and the density $R$}\label{sec:measures}

\begin{definition}[Full and on-critical measures]\label{def:measures}
Let $dm_\infty$ be the reference (archimedean) measure of the spectral triple. The full and
on-critical spectral measures are the absolutely continuous measures
\[
\dmfull=|\zeta(\half+is)|^2\,dm_\infty,\qquad
\dmon=|\zon(\half+is)|^2\,dm_\infty,
\]
where $\zon$ is defined by the Hadamard product of $\zeta$ with every nontrivial zero
$\rho=\half+\delta+i\gamma$ replaced by its on-critical projection $\half+i\gamma$. The density
of $\dmfull$ relative to $\dmon$ is $R(s)=|\zeta(\half+is)/\zon(\half+is)|^2$, and
\[
\dnu=\dmfull-\dmon=(R-1)\,\dmon .
\]
\end{definition}

\section{The dominance theorem}\label{sec:dominance}

\begin{theorem}[Unconditional dominance of $\zeta$ over $\zon$]\label{thm:dominance}
For all $s\in\mathbb R$, $R(s)\ge1$, with equality for all $s$ if and only if $\RH$ holds.
\end{theorem}
\begin{proof}
The nontrivial zeros are symmetric under conjugation and $s\mapsto1-s$, so each off-critical
zero $\rho_0=\half+\delta_0+i\gamma_0$ ($\delta_0\ne0$) belongs to a quadruple
$\{\rho_0,\bar\rho_0,1-\rho_0,1-\bar\rho_0\}$; the on-critical replacement moves all four to
$\half\pm i\gamma_0$ (each with multiplicity two). On the critical line $s=\half+it$, the ratio
$\zeta/\zon$ is the product over off-critical quadruples of
\[
\frac{(\half+it-\rho_0)(\half+it-\bar\rho_0)(\half+it-(1-\rho_0))(\half+it-(1-\bar\rho_0))}
{(\half+it-(\half+i\gamma_0))^2(\half+it-(\half-i\gamma_0))^2}.
\]
Computing moduli, $|\half+it-\rho_0|^2=\delta_0^2+(t-\gamma_0)^2$, and likewise the four
numerator factors give $\delta_0^2+(t\mp\gamma_0)^2$ (twice each), while the denominators give
$(t-\gamma_0)^2$ and $(t+\gamma_0)^2$ (squared). Hence the squared modulus of the
single-quadruple contribution is
\[
\frac{(\delta_0^2+(t-\gamma_0)^2)^2(\delta_0^2+(t+\gamma_0)^2)^2}{(t-\gamma_0)^4(t+\gamma_0)^4}
\ge1,
\]
since $\delta_0^2+u^2\ge u^2$ for every $u$, with equality iff $\delta_0=0$. The product over
quadruples converges absolutely (the off-critical zeros, if any, obey the global density bound
$N(T)=O(T\log T)$, and all factors are $\ge1$), giving $R(s)\ge1$, with equality for all $s$ iff
every $\delta_0=0$, i.e.\ $\RH$.
\end{proof}

\begin{corollary}[$\dnu$ is a non-negative invariant for $\RH$]\label{cor:dnu}
Unconditionally $\dnu=(R-1)\,\dmon\ge0$. Moreover $\dnu=0\iff\RH$, and $\dnu$ is even
($\dnu(s)=\dnu(-s)$, from $\xi(s)=\xi(1-s)$) and finite, with
\[
\int_{\mathbb R}\dnu=\sum_{j}\frac{4\delta_j^2}{\gamma_j^2+\delta_j^2}+O\!\Big(\sum_j
\frac{\delta_j^4}{\gamma_j^4}\Big)<\infty,
\]
the sum over off-critical zeros, convergent because $\sum_j\delta_j^2/\gamma_j^2<\infty$ by the
zero-counting estimate.
\end{corollary}
\begin{proof}
Non-negativity and the equivalence are Theorem~\ref{thm:dominance}. Parity is the functional
equation. For finiteness, each quadruple contributes $(R-1)$ of order $\delta_j^2/\gamma_j^2$,
summable by $N(T)=O(T\log T)$ and $\sum_j\delta_j^2/\gamma_j^2<\infty$.
\end{proof}

\begin{remark}[The curve case]\label{rem:curve}
For a smooth projective curve over a finite field, the analogous measure vanishes identically:
every zero of the zeta function already lies on the critical circle, so the product analogous
to $R$ has no nontrivial factors. There $\dnu=0$ is \emph{known} from Weil's proof; for $\zeta$
we know only $\dnu\ge0$. The whole difficulty is the gap between these two statements, made
precise in \S\ref{sec:weil}--\S\ref{sec:barrier}.
\end{remark}

\section{The synthesis: eight equivalent criteria}\label{sec:eight}

Several criteria that arise independently --- from a Nevanlinna transform, a Jensen excess on a
disk model, a relative entropy, and the spectral trace --- all reduce to $\dnu=0$.

\begin{theorem}[Eight equivalent criteria]\label{thm:eight}
The following are pairwise equivalent.
\begin{enumerate}[\upshape(1)]
\item \textup{(RH)} every nontrivial zero of $\zeta$ has $\Re s=\half$;
\item \textup{(Zero measure)} $\dnu=0$;
\item \textup{(Nevanlinna)} $\mathcal G(z)=\int_{\mathbb R}\frac{s+z}{s-z}\,\dnu(s)$ vanishes
identically on $\{\Im z>0\}$;
\item \textup{(Trace)} $\Tlam=\int_{\mathbb R}\Wlam\,\dnu=0$ for every $\lambda>0$, where
$\Wlam>0$ is the Abel kernel;
\item \textup{(Jensen excess)} $\delta=\sum_{\Re\rho>1/2}\log(|\rho|/|\rho-1|)=0$;
\item \textup{(Variational)} the squared distance $\widetilde\delta_\lambda^2$ from
$|\zeta|^2dm_\infty$ to the on-critical variety vanishes for every $\lambda$;
\item \textup{(Relative entropy)} $\mathrm{KL}(\dmfull\,\|\,\dmon)=0$;
\item \textup{(CCM distance)} $d_{\mathrm{CCM}}=\sup_{\lambda>0}\Tlam=0$.
\end{enumerate}
The equivalences $(1)\Leftrightarrow(2)\Leftrightarrow(3)\Leftrightarrow(5)\Leftrightarrow(7)$
are unconditional. The equivalences involving the trace, $(2)\Leftrightarrow(4)\Leftrightarrow
(6)\Leftrightarrow(8)$, hold under the completeness hypothesis of \S\ref{sec:trace}
\textnormal{(}the totality of $\{\Wlam\}$ in $L^2(\dnu)$\textnormal{)}; the forward directions
$(2)\Rightarrow(4),(6),(8)$ are unconditional.
\end{theorem}
\begin{proof}
$(1)\Leftrightarrow(2)$ is Corollary~\ref{cor:dnu}. $(2)\Leftrightarrow(3)$: the Nevanlinna
kernels $\{(s+z)/(s-z):\Im z>0\}$ are total in $L^2(\dnu)$, so $\mathcal G\equiv0\iff\dnu=0$.
$(1)\Leftrightarrow(5)$: by Jensen's formula on the disk model $F(z)=\xi(z/(z-1))$, the excess
$\delta\ge0$ with equality iff $F$ is zero-free in $|z|<1$, iff $\RH$. $(2)\Leftrightarrow(7)$:
$\mathrm{KL}(\dmfull\|\dmon)=\int\log R\,\dmon$, and $\log x\ge(x-1)-\half(x-1)^2$ with $R\ge1$
gives $\mathrm{KL}=0\iff R=1$ a.e.\ $\iff\dnu=0$. For the trace equivalences: $\dnu\ge0$ and
$\Wlam>0$ give $\Tlam\ge0$, so $\dnu=0\Rightarrow\Tlam=0\ \forall\lambda$ (forward,
unconditional); the converse $\Tlam=0\ \forall\lambda\Rightarrow\dnu=0$ is exactly the
completeness hypothesis of \S\ref{sec:trace}. $(2)\Leftrightarrow(6)$ is the same statement,
since $\widetilde\delta_\lambda^2=\Tlam$ up to a positive constant, and $(4)\Leftrightarrow(8)$
because $d_{\mathrm{CCM}}=\sup_\lambda\Tlam$ vanishes iff every $\Tlam$ does.
\end{proof}

\section{The trace, its lower bound, and the completeness hypothesis}\label{sec:trace}

\begin{proposition}[Lower bound under $\neg\RH$]\label{prop:lower}
Suppose $\neg\RH$, with off-critical zeros $\rho_j=\half+\delta_j+i\gamma_j$. For each
$\lambda>0$ and each $j$ with $|\gamma_j|\le\lambda$,
\[
\Tlam\ge2\sum_{|\gamma_j|\le\lambda}\delta_j^2\,K_\lambda(\gamma_j),\qquad
K_\lambda(\gamma_j)=\Wlam(\gamma_j)\,\dmon(\{\gamma_j\})>0
\]
for $\lambda$ large enough; hence the right-hand side is strictly positive. The trace detects an
off-critical zero to second order in its displacement $\delta_j$.
\end{proposition}
\begin{proof}
By Theorem~\ref{thm:dominance}, near $\gamma_j$ the density satisfies $R(t)-1\ge(\delta_j^2+
(t-\gamma_j)^2)^2/(t-\gamma_j)^4-1\ge2\delta_j^2/(t-\gamma_j)^2+O(\delta_j^4)$; integrating
$\Wlam\,(R-1)\,\dmon$ over a unit window at $\gamma_j$ and summing the (non-negative)
contributions gives the bound.
\end{proof}

The decisive point is the reverse implication. Since $\dnu\ge0$ and $\Wlam>0$, the trace is a
\emph{linear} functional of a non-negative measure tested against a positive kernel; there is no
sign cancellation, and the only question is whether the kernels $\{\Wlam\}$ \emph{see all of}
$\dnu$.

\begin{definition}[Completeness hypothesis]\label{def:completeness}
\textbf{(C)}\ The Abel kernels $\{\Wlam:\lambda>0\}$ are total in $L^2(\dnu)$ for every
non-negative even finite measure $\dnu$ of the form $(R-1)\dmon$.
\end{definition}

\begin{theorem}[Trace criterion under (C)]\label{thm:trace-crit}
Under hypothesis \textnormal{(C)}, $\Tlam=0$ for all $\lambda>0$ if and only if $\RH$.
Unconditionally, $\RH\Rightarrow\Tlam=0\ \forall\lambda$, and a single off-critical zero forces
some $\Tlam>0$ \textnormal{(}Proposition~\ref{prop:lower}\textnormal{)}.
\end{theorem}
\begin{proof}
$\RH\Rightarrow\dnu=0\Rightarrow\Tlam=0$ for all $\lambda$ (Corollary~\ref{cor:dnu}). Under (C),
$\Tlam=0\ \forall\lambda$ means $\dnu\perp\Wlam$ for all $\lambda$, hence $\dnu=0$ by totality,
hence $\RH$.
\end{proof}

\begin{remark}[What is and is not used]\label{rem:notused}
We emphasize that Theorem~\ref{thm:trace-crit} uses \emph{only} $\dnu\ge0$ (the dominance
theorem) and the scalar positivity $\Wlam>0$. It does \emph{not} use, and we do not assert, any
positivity of $\Wlam$ as a quadratic form (a kernel positivity); such a pointwise-kernel claim
is false in general and plays no role here. Hypothesis (C) is a genuine analytic input --- a
completeness of the Abel system in $L^2(\dnu)$ --- and is the honest residue of the trace
criterion. It is equivalent, via the orthogonal-polynomial realization of $\dnu$, to the
moment-rigidity criterion proved in \cite{MomentRigidity}, where $\RH$ is shown equivalent to
the constancy of a sequence of Christoffel-type functionals; (C) is the statement that the trace
family recovers those functionals.
\end{remark}

\section{The function-field mirror: Weil's three pillars}\label{sec:weil}

The reason $\dnu$ vanishes in the function-field case is structural and instructive.

\begin{theorem}[Function-field trace vanishes]\label{thm:ff}
For the zeta function $Z_C$ of a smooth projective curve $C$ over a finite field, the analogue
$\dnu^C$ of $\dnu$ is identically zero.
\end{theorem}

The vanishing rests on three pillars, of which $\zeta$ has analogues of the first two but not
the third:
\begin{enumerate}[(P1)]
\item a finite-dimensional cohomology $H^1(C)$ carrying the zeros as Frobenius eigenvalues
(\emph{for $\zeta$}: a candidate spectral realization exists --- the Jacobi operator of the
spectral triple --- but it is infinite-dimensional and its self-adjointness is the analogue of,
not a substitute for, the Frobenius);
\item a Lefschetz/explicit-formula trace relating the eigenvalues to point counts
(\emph{for $\zeta$}: the Riemann--Weil explicit formula plays this role);
\item a \textbf{polarization}: a positive intersection form on $H^1(C)$ (the Hodge index /
Riemann form) that forces the Frobenius eigenvalues onto the critical circle
(\emph{for $\zeta$}: absent --- this is exactly the missing positivity).
\end{enumerate}

\begin{remark}
The measure $\dnu$ is the analytic shadow of the failure of: in the curve case the
polarization forces $R\equiv1$ (all zeros on the circle, $\dnu^C=0$); for $\zeta$ the Hadamard
product gives only $R\ge1$, the one-sided remnant of a polarization that is not known to be
two-sided. The trace criterion (Theorem~\ref{thm:trace-crit}) and the moment rigidity of
\cite{MomentRigidity} are two ways of asking for the missing pillar.
\end{remark}

\section{The Hadamard barrier}\label{sec:barrier}

\begin{definition}[Hadamard barrier]\label{def:barrier}
The \emph{Hadamard barrier} is the gap between the two facts: (a) the Hadamard product of $\xi$
yields, unconditionally, the one-sided bound $R\ge1$, i.e.\ $\dnu\ge0$; (b) forcing $\dnu=0$
requires either the completeness (C) of the Abel system in $L^2(\dnu)$ or, equivalently, a
positive intersection form on the spectral realization. The Hadamard product delivers (a)
but not (b).
\end{definition}

\begin{proposition}[The barrier is genuine]\label{prop:barrier}
The implication $\dnu\ge0\Rightarrow\dnu=0$ is false for general non-negative even finite
measures; it holds for the specific $\dnu=(R-1)\dmon$ if and only if $\RH$. Hence no purely
soft argument from $\dnu\ge0$ can close the gap: the additional input is exactly of the strength
of $\RH$, and is supplied by neither the Hadamard product nor the explicit formula alone.
\end{proposition}
\begin{proof}
A non-negative even finite measure need not vanish, so $\dnu\ge0$ alone is insufficient. By
Corollary~\ref{cor:dnu}, $\dnu=(R-1)\dmon=0\iff\RH$. Thus any argument deducing $\dnu=0$ supplies
$\RH$, so it cannot be soft; in particular it is not a consequence of the structural facts
(dominance, explicit formula, self-adjointness of the Jacobi operator) that hold unconditionally.
\end{proof}

\section{Conclusion}

The difference of the two spectral measures of the zeta spectral triple is a single, clean,
non-negative invariant: $\dnu\ge0$ unconditionally (Theorem~\ref{thm:dominance}), with $\dnu=0$
exactly the Riemann Hypothesis (Corollary~\ref{cor:dnu}). It unifies a family of criteria
(Theorem~\ref{thm:eight}) and is detected to second order by a positive trace family
(Proposition~\ref{prop:lower}). The criterion's hard direction is honestly isolated: it is the
completeness (C) of the Abel kernels in $L^2(\dnu)$ (Theorem~\ref{thm:trace-crit},
Remark~\ref{rem:notused}), equivalent to the moment-rigidity criterion of \cite{MomentRigidity},
and we use no pointwise kernel positivity. The function-field mirror (\S\ref{sec:weil}) names
the missing ingredient as Weil's polarization pillar, and the Hadamard barrier
(\S\ref{sec:barrier}) proves that no soft argument from $\dnu\ge0$ can supply it: the input
needed is exactly of the strength of $\RH$. The paper neither assumes nor proves the Riemann
Hypothesis; it provides a sign-definite invariant whose vanishing is $\RH$ and a precise map of
what its vanishing demands.

\begin{thebibliography}{99}
\bibitem{MomentRigidity} D.~A.~Trejo Pizzo, \emph{A moment-rigidity criterion for the Riemann
Hypothesis via Meixner--Pollaczek polynomials}, preprint, 2026.
\bibitem{CCM} A.~Connes, C.~Consani, H.~Moscovici, \emph{Zeta spectral triples}, preprint,
arXiv:2511.22755, 2025.
\bibitem{BaezDuarte} L.~B\'aez-Duarte, \emph{A strengthening of the Nyman--Beurling criterion for
the Riemann hypothesis}, Atti Accad.\ Naz.\ Lincei \textbf{14} (2003), 5--11.
\bibitem{Weil} A.~Weil, \emph{Sur les ``formules explicites'' de la th\'eorie des nombres
premiers}, Comm.\ S\'em.\ Math.\ Univ.\ Lund (1952), 252--265.
\bibitem{KS} I.~S.~Kac, M.~G.~Kre\u\i n, \emph{R-functions --- analytic functions mapping the
upper half-plane into itself}, Amer.\ Math.\ Soc.\ Transl.\ \textbf{103} (1974), 1--18.
\bibitem{bgConnes99} A.~Connes, \emph{Trace formula in noncommutative geometry and the zeros of the Riemann zeta function}, Selecta Math.\ (N.S.) \textbf{5} (1999), 29--106.
\bibitem{bgCC21} A.~Connes and C.~Consani, \emph{Weil positivity and trace formula, the archimedean place}, Selecta Math.\ (N.S.) \textbf{27} (2021), no.~77.
\end{thebibliography}

\end{document}
